\chapter{Puissances et racines carrées}
\label{corriges:04}
\section*{Exercice 4.1 -- Les exposants ont un sens}
$2^{-3}=1/8$.
\[(3^2)^3/3^4=3^{6-4}=3^2=9.\]
Le quotient scientifique vaut $(8/2)\times10^{7-(-3)}=4\times10^{10}$. Soustraire un exposant négatif revient à ajouter sa valeur opposée.

\section*{Exercice 4.2 -- Racines et géométrie}
Le côté vaut $\sqrt{98}=\sqrt{49\times2}=7\sqrt2\unite{cm}\approx9{,}8995\unite{cm}$. Au millimètre près : $9{,}9\unite{cm}$, soit $99\unite{mm}$.
$\sqrt{(-8)^2}=\sqrt{64}=8=\abs{-8}$, et non $-8$.

\section*{Exercice 4.3 -- Doublements}
$4$ heures contiennent $8$ périodes de $30$ minutes. Le modèle donne $150\times2^8=150\times256=38400$ cellules. L'exposant compte les doublements, pas le nombre d'heures. La validité biologique du modèle reste limitée par ses hypothèses.

\section*{Exercice 4.4 -- Réfuter puis expliquer}
Avec $a=9,b=16$, le premier membre vaut $5$ et le second $7$. Pour le produit, $\sqrt a\sqrt b$ est positif ou nul et son carré vaut $(\sqrt a)^2(\sqrt b)^2=ab$. Par unicité de la racine carrée positive ou nulle, il est égal à $\sqrt{ab}$.
